\documentclass[11pt]{article}

\usepackage[letterpaper,margin=0.85in]{geometry}
\usepackage[T1]{fontenc}
\usepackage[utf8]{inputenc}
\usepackage{lmodern}
\usepackage{amsmath,amssymb}
\usepackage{booktabs}
\usepackage{enumitem}
\usepackage{xcolor}
\usepackage{fancyhdr}
\usepackage{hyperref}
\usepackage{listings}

\definecolor{GTNavy}{HTML}{003057}
\definecolor{GTGold}{HTML}{B3A369}
\definecolor{GTBlue}{HTML}{0076A8}
\definecolor{GTLight}{HTML}{F3F0E8}
\definecolor{GTGray}{HTML}{54585A}

\hypersetup{colorlinks=true,linkcolor=GTBlue,urlcolor=GTBlue}
\pagestyle{fancy}
\fancyhf{}
\lhead{ISyE 4031: Regression and Forecasting}
\rhead{Public sample}
\cfoot{\thepage}
\setlength{\headheight}{14pt}
\setlist[itemize]{leftmargin=1.5em,itemsep=0.25em}
\setlist[enumerate]{leftmargin=1.8em,itemsep=0.35em}
\setlength{\parindent}{0pt}
\setlength{\parskip}{0.55em}

\lstset{
  basicstyle=\ttfamily\small,
  backgroundcolor=\color{GTLight},
  frame=single,
  rulecolor=\color{GTGold},
  breaklines=true,
  showstringspaces=false,
  columns=fullflexible
}

\newcommand{\points}[1]{\nobreak\hfill\mbox{\textcolor{GTBlue}{\textbf{[#1 points]}}}}

\begin{document}

{\color{GTGold}\rule{\textwidth}{3pt}}
\vspace{0.5em}

{\LARGE\bfseries\color{GTNavy}Homework 1: Statistical Review, Forecast Baselines, and Agent Audit}

\vspace{0.4em}
\begin{tabular}{@{}ll@{}}
\textbf{Course:} & ISyE 4031 \\
\textbf{Version:} & Public sample \\
\textbf{Work mode:} & Individual \\
\textbf{Format:} & Reproducible project folder or ZIP archive
\end{tabular}

\section*{Purpose}

This assignment asks you to use statistical reasoning and an AI coding agent together. The agent may propose code, explanations, and tests, but you remain responsible for the estimand, assumptions, calculations, information cutoff, and conclusions.

Your files must allow another person to reproduce every reported result without editing your code.

\section*{Required deliverables}

Create a project folder containing the following files at the top level:
\begin{itemize}
  \item \texttt{analysis.py} or a clearly organized notebook that produces every numerical result;
  \item \texttt{test\_analysis.py} with the required automated checks;
  \item \texttt{report.pdf}, no more than five pages, answering all written questions;
  \item \texttt{AI\_USAGE.md}, using the evidence requirements in Part 4;
  \item a short \texttt{README.md} with exact commands for reproducing the analysis and tests.
\end{itemize}

If you share the project as a ZIP archive, name it \texttt{hw1\_lastname\_firstname.zip}. Do not include a Git repository, virtual environment, package cache, credentials, or unrelated files.

Run the following before considering the work complete:
\begin{lstlisting}[language=bash]
python analysis.py
pytest -q
\end{lstlisting}

Then create the optional archive. If you used a notebook, substitute its filename for \texttt{analysis.py}.
\begin{lstlisting}[language=bash]
zip hw1_lastname_firstname.zip analysis.py test_analysis.py \
    report.pdf AI_USAGE.md README.md
\end{lstlisting}

\section*{Part 1: Descriptive summaries and uncertainty}

The following course-provided practice data represent 20 consecutive daily maximum temperatures in degrees Fahrenheit. Treat the first 14 observations as the legal information set and the final 6 observations as a held-out evaluation period.

\begin{lstlisting}[language=Python]
temps = [86, 88, 91, 90, 93, 95, 94, 92, 89, 87,
         88, 90, 92, 94, 96, 95, 93, 91, 90, 88]
train = temps[:14]
test = temps[14:]
\end{lstlisting}

Do not use the held-out values when estimating a model or interval.

\begin{enumerate}[label=\textbf{1\alph*.}]
  \item Plot the 20 temperatures in their observed order. Mark the cutoff after day 14, distinguish training from held-out observations, label both axes, and write a one-sentence caption explaining what information is legal when the forecasts are issued. \points{4}
  \item State the unit of analysis, response, information cutoff, and intended forecast target. \points{4}
  \item Using only \texttt{train}, calculate the mean, median, range, sample variance, sample standard deviation, and interquartile range. Show enough work to verify the denominator used for the variance. \points{6}
  \item Assuming the training observations are independent draws from an approximately normal population, construct a 95\% t confidence interval for the population mean. State the formula and interpret the interval as a repeated-sampling procedure. \points{5}
  \item Under the same assumptions, construct a 95\% prediction interval for one new daily maximum temperature. Explain why this interval is wider than the interval in 1d. \points{6}
\end{enumerate}

\section*{Part 2: A transparent forecast baseline}

Use the same fixed training set to create two six-step forecasts:
\begin{itemize}
  \item \textbf{Mean baseline:} forecast every held-out day using the training mean.
  \item \textbf{Last-value baseline:} forecast every held-out day using the final observed training value.
\end{itemize}

\begin{enumerate}[label=\textbf{2\alph*.}]
  \item Report the six forecasts from each baseline and overlay both forecast paths on the ordered-data plot from Part 1. Explain why neither forecast may change after the held-out period begins. \points{5}
  \item Calculate mean absolute error and root mean squared error for both baselines on \texttt{test}. Identify the better baseline under each metric. \points{8}
  \item Explain one reason these six errors are not six independent replications of model performance. \points{4}
  \item Choose one baseline as the required benchmark for a future extension. Defend the choice using simplicity, performance, and reproducibility. \points{5}
\end{enumerate}

\section*{Part 3: One-sample inference}

A melting-point test of $n=10$ binder samples produced $\bar x=154.2^\circ$F. Assume the population is normal with known standard deviation $\sigma=1.5^\circ$F.

\begin{enumerate}[label=\textbf{3\alph*.}]
  \item Test $H_0:\mu=155$ against $H_a:\mu\neq155$ at $\alpha=0.01$. Show the test statistic, rejection rule, decision, and conclusion in context. \points{8}
  \item Calculate and interpret the two-sided p-value. Do not describe it as the probability that $H_0$ is true. \points{5}
  \item Calculate the probability of a Type II error when the true mean is $150^\circ$F. State the acceptance region used in the calculation. \points{7}
  \item Verify all three results in Python and compare them with your independent calculation. \points{4}
\end{enumerate}

\section*{Part 4: Audit an agent-generated analysis}

An AI agent proposes the following code for Part 1 and Part 2:

\begin{lstlisting}[language=Python]
import numpy as np
from scipy import stats

x = np.array(temps, dtype=float)
train = x[:14]
test = x[14:]

center = x.mean()
spread = train.std()
se = spread / len(train)
mean_interval = stats.norm.interval(.95, loc=center, scale=se)
future_interval = mean_interval
forecast = np.repeat(x[-1], len(test))
\end{lstlisting}

\begin{enumerate}[label=\textbf{4\alph*.}]
  \item Identify at least five distinct statistical, forecasting, or reproducibility problems. For each problem, explain the consequence rather than merely naming the faulty line. \points{10}
  \item Replace the code with a correct implementation. Keep the two transparent baselines separate from any later extension. \points{6}
  \item Add automated tests demonstrating all four properties below: \points{8}
    \begin{itemize}
      \item changing held-out values cannot change training summaries or forecasts;
      \item the sample standard deviation uses $n-1$ in the denominator;
      \item the prediction interval is wider than the confidence interval for the mean;
      \item repeated runs produce identical results.
    \end{itemize}
\end{enumerate}

\section*{AI Usage Statement}

Use at least one designated course agent for Part 4. Your \texttt{AI\_USAGE.md} must include:
\begin{itemize}
  \item the tool and model used;
  \item one representative prompt or prompt template;
  \item the task delegated to the agent;
  \item the files or functions the agent read or changed;
  \item the tests the agent proposed or ran;
  \item one suggestion you accepted, revised, or rejected, with your statistical reason;
  \item one unresolved limitation or risk.
\end{itemize}

Do not include private chain-of-thought or unrelated account history. A concise project-specific audit record is sufficient. \points{5}

\section*{Rubric summary}

\begin{center}
\begin{tabular}{lr}
\toprule
\textbf{Component} & \textbf{Points} \\
\midrule
Descriptive summaries and uncertainty & 25 \\
Forecast baselines and evaluation & 22 \\
One-sample inference & 24 \\
Agent audit and automated tests & 24 \\
AI Usage Statement & 5 \\
\midrule
\textbf{Total} & \textbf{100} \\
\bottomrule
\end{tabular}
\end{center}

\section*{Reproducibility checks}

Before considering the work complete, confirm that:
\begin{itemize}
  \item the ZIP file opens and contains the five required files at its top level;
  \item all reported numbers can be regenerated from the project files;
  \item no held-out observation enters training, interval estimation, or baseline construction;
  \item every assumption and interpretation is written in your own defensible language;
  \item \texttt{pytest -q} passes in the course Codespace;
  \item no credentials, tokens, private data, or unrelated files are included.
\end{itemize}

\end{document}
